\documentclass[a4paper]{article}

\usepackage{graphicx}
\usepackage{amsmath,amssymb}
\usepackage{minted}
\usepackage{multirow}
\usepackage{float}
\usepackage{todonotes}
\usepackage{forest}
\usepackage{xstring}
\usepackage{xcolor}
\usepackage[most]{tcolorbox}
\usepackage{xparse}
\usepackage{natbib}

\usetikzlibrary{graphs,graphdrawing}
\usegdlibrary{layered}

% for external graphviz
\input{/home/donkarlo/Dropbox/repo/nd_ai_project/src/nd_ai/knoweldge_representation/language/formal/computer/programming/tex/latex/code_clips/external_graphviz.tex}

% for tags
\input{/home/donkarlo/Dropbox/repo/nd_ai_project/src/nd_ai/knoweldge_representation/language/formal/computer/programming/tex/latex/parts/control_sequence/kind/semtag/semtag.tex}

% for links
\input{/home/donkarlo/Dropbox/repo/nd_ai_project/src/nd_ai/knoweldge_representation/language/formal/computer/programming/tex/latex/code_clips/seehere.tex}

\title{Depolarization}
\date{}
\author{Mohammad Rahmani}

\begin{document}
    \maketitle
    In biology, depolarization or hypopolarization is a change within a cell, during which the cell undergoes a shift in electric charge distribution, resulting in less negative charge inside the cell compared to the outside \seeherefooter{https://en.wikipedia.org/wiki/Depolarization}.
    
    An action potential is a rapid, all-or-none change in membrane voltage across an excitable cell membrane. Depolarization is mainly caused by sodium influx, while repolarization is mainly caused by potassium efflux. Ion gradients and the Na/K-ATPase pump help restore and maintain the resting membrane potential after excitation \cite{grider-2023-physiology-action-potential}.
    
    \bibliographystyle{plainnat}
    \bibliography{/home/donkarlo/Dropbox/repo/nd_shared_research/bibliography/bibliography.bib}
\end{document}